\begin{table}[htbp]\centering\footnotesize
\caption{Los catorce ramos que más se mueven en pesos, línea P (mdp)}\label{cua:ramos}
\begin{tabular}{llrrrr}
\toprule
Ramo & Denominación & 2026 & 2027 & Dif. & Var. real (\%) \\
\midrule
18 & Energía & 267.439,1 & 85.634,2 & $-$181.804,9 & $-$69,2 \\
19 & Aportaciones a Seguridad Social & 1.541.518,7 & 1.661.007,9 & 119.489,2 & +3,6 \\
33 & Aportaciones Federales para Entidades… & 1.041.892,9 & 1.136.351,6 & 94.458,7 & +4,9 \\
28 & Participaciones a Entidades Federativa… & 1.456.045,9 & 1.547.607,9 & 91.562,0 & +2,2 \\
11 & Educación Pública & 513.015,6 & 576.631,9 & 63.616,3 & +8,1 \\
34 & Erogaciones para los Programas de Apoy… & 35.553,4 & 65.906,1 & 30.352,7 & +78,2 \\
20 & Bienestar & 674.510,0 & 697.106,3 & 22.596,3 & $-$0,6 \\
56 & Servicios de Salud del Instituto Mexic… & 172.492,4 & 193.987,2 & 21.494,8 & +8,1 \\
22 & Instituto Nacional Electoral & 22.837,2 & 42.274,5 & 19.437,3 & +78,0 \\
13 & Marina & 65.926,8 & 81.223,4 & 15.296,6 & +18,5 \\
07 & Defensa Nacional & 170.753,1 & 185.706,1 & 14.953,0 & +4,6 \\
03 & Poder Judicial & 85.960,2 & 72.264,7 & $-$13.695,5 & $-$19,2 \\
08 & Agricultura y Desarrollo Rural & 75.195,5 & 88.758,8 & 13.563,3 & +13,5 \\
09 & Infraestructura, Comunicaciones y Tran… & 153.539,3 & 167.074,8 & 13.535,5 & +4,6 \\
\bottomrule
\end{tabular}
\\[2pt]\begin{minipage}{0.92\textwidth}\scriptsize Fuente: Anexo 1 de los proyectos de decreto de PEF 2026 y 2027, \textbf{por clave de ramo, no por nombre}. Deflactor del PIB de 2027. \textbf{Un ramo no es una función}: una caída de ramo no es una caída de la política pública que su nombre sugiere.\end{minipage}
\end{table}
